% Generated by Claude Opus 5 (Anthropic)
% The four-quadrant model figure (drawn in Miro, exported as one PDF frame).
%
% Overleaf steps:
%   1. In Miro select the frame (click its title bar) -> ... -> Save as PDF.
%   2. Upload it into the Overleaf project as  figures/model_panel.pdf
%      (a PDF keeps the text vector; PNG only if the PDF export is unavailable,
%      then export at 4x and name it figures/model_panel.png).
%   3. \usepackage{graphicx} in the preamble, then \input this file, or just
%      paste the block below into the paper.
%   4. Refer to it as Fig.~\ref{fig:model}.
%
% figure* spans both columns; the frame is ~2.3:1, so width=\textwidth gives
% about 3.1 in of height on IEEEtran -- roughly a third of the page.

\begin{figure*}[t]
  \centering
  \includegraphics[width=\textwidth]{figures/model_panel}
  \caption{What FISH extracts from an unmodified ROS~2 + GPU stack, on one
    Autoware run.  \textbf{(a)}~the containment hierarchy from the host down to
    callbacks, with the namespaces that group executors and nodes;
    \textbf{(b)}~the same fragment across layers: the per-lidar filters, the
    concatenation node whose two publishing callbacks form a join, the
    sample-and-hold (state) reads of the distortion corrector, and the
    CenterPoint detector that submits GPU work;
    \textbf{(c)}~the events each edge is decided from, and the edge each
    pattern yields;
    \textbf{(d)}~for the GPU-submitting callback, the typed DAG of every
    execution (CPU calls, kernels on the SMs, transfers on the copy engines)
    and the conditional graph that summarises them: the part common to all
    executions plus the branches, each with the share of executions that take
    it.}
  \label{fig:model}
\end{figure*}
